%BehaviorSemantics.tex

\ksec{Behaviors (KerML 8.4.4.7.1)}

\texttt{Behaviors}\index{behavior} are \texttt{Classes} which implicitly specialize \texttt{Performances::Performance} which is an \texttt{Occur\-rence} disjoint from \texttt{Object}.  Because structures specialize \texttt{Object}, behaviors and structures must be disjoint.  A behavior exists in time because it is an \texttt{Occur\-rence}.  

Behaviors may have parameters\index{parameter} which are features having direction \texttt{in}\index{in}, \texttt{out}\index{out}, or \texttt{inout}\index{inout}.  Curiously, there is no definition as to the meaning of directed features other than input, output, or both (see \S\ref{sec:variables}).

\ecaption{behavior}
\begin{lstlisting}
// Specializes Performances::Performance by default.
behavior TakePicture {
  in scene : Scene;
  out picture : Picture;
}
\end{lstlisting}

Parameters are ordered in the lexical order they are declared in the body of a behavior. They may appear at any
location within the body.

If a behavior (subclassifier) specializes another behavior (superclassifier), then each of the owned parameters
of the subclassifier behavior must, in order, redefine the parameter at the same position of each of the superclassifier
behaviors. The redefining parameters shall have the same direction as the redefined parameters.

Here, behavior \texttt{C} specializes both \texttt{A} and \texttt{B}:
\ecaption{behavior specializes}
\begin{lstlisting}
behavior A { in a1; out a2; }
behavior B { in b1; out b2; }
behavior C specializes A, B {
  in c1 redefines a1, b1;
  out c2 redefines a2, b2;
}
\end{lstlisting}

Because subclassifier behavior parameters implicitly redefine their superclassifers' parameters in order, this is equivalent:
\ecaption{behavior specializes equivalent}
\begin{lstlisting}
behavior C specializes A, B {
  in c1;
  out c2;
}
\end{lstlisting}

Behaviors having only parameters don't actually do anything.  For that other features are needed.

\ecaption{TakePicture}
\begin{lstlisting}
behavior Focus { in scene : Scene; out image : Image; }
behavior Shoot { in image : Image; out picture : Picture; }
behavior TakePicture {
  in scene : Scene;
  out picture : Picture;
  binding focus.scene = scene;
  composite step focus : Focus;
  succession focus then shoot;
  composite flow focus.image to shoot.image;
  composite step shoot : Shoot;
  binding picture = shoot.picture;
}
\end{lstlisting}

\pn By \S\ref{par:cnp}, composite steps \texttt{focus} and \texttt{shoot} occur during \texttt{TakePicture}:\\
\texttt{during(focus,TakePicture)}~~\texttt{during(shoot,TakePictrue)}\\
By \ref{eq:df-succession} and \texttt{succession focus then shoot}:\\
\texttt{precedes(focus,shoot)}\\
By \ref{eq:df-bind} input and output parameters are forced to be identical.  There does not seem to be direction restrictions for parameter binding.\\
By \S\ref{sec:flows} a value is transferred from output \texttt{focus.image} to input \texttt{shoot.image} by a flow.

\ksec{Steps (KerML 8.4.4.7.2)}

\pn A \texttt{Step}\index{step} is a \texttt{feature} typed by a \texttt{Behavior}.  Anything can have a step--not just behaviors.

\pn Steps are most useful when ordered by successions or succession flows.  Otherwise, the occurrences of steps can overlap.

\pn Parameter features of steps can be connected with binding connectors.

Steps declared in the body of a behavior are the owned steps of the featuring behavior.
\ecaption{step}
\begin{lstlisting}
behavior Focus { in scene : Scene; out image : Image; }
behavior Shoot { in image : Image; out picture : Picture; }
behavior TakePicture {
  in scene : Scene;
  out picture : Picture;
  composite step focus : Focus;
  composite step shoot : Shoot;
}
\end{lstlisting}

\pn \texttt{Performance} has both \texttt{enclosedPerformances} and \texttt{subperformances} which is \texttt{composite}.  They seem redundant because time-enclosed performances have lifetimes within their parent performance, which is the same meaning as ``composite".

\ecaption{enclosedPerformances}
\begin{lstlisting}
		step enclosedPerformances: Performance[0..*] subsets performances, timeEnclosedOccurrences
			intersects performances, timeEnclosedOccurrences {
			doc
			/*
			 * timeEnclosedOccurrences of this Performance that are also Performances.
			 */
		}
\end{lstlisting}

\ecaption{subperformances}
\begin{lstlisting}
composite step subperformances: Performance[0..*] subsets enclosedPerformances, suboccurrences
	intersects enclosedPerformances, suboccurrences {
  feature redefines this default (that as Performance).this {
    doc
		/*
		 * The default "this" context of a subperformance is the same as that of its owning Performance.
		 * This means that the context for any Performance that is in a composition tree rooted
		 * in a Performance that is not itself owned by an Object is the root Performance. If the
		 * root Performance is an ownedPerformance of an Object, then that Object is the context.
		 */
	}
	feature redefines thisPerformance default (that as Performance).thisPerformance;
}		
\end{lstlisting}
\comment{What?}

\pn A \texttt{Step} whose \texttt{owningType} is a \texttt{Behavior} or another \texttt{Step} must directly or indirectly subset
 \texttt{Performan\-ces::Performance::enclosedPerformances}.

\pn A composite \texttt{Step} whose \texttt{owningType} is a \texttt{Structure} or a \texttt{Feature} typed by a \texttt{Structure} must directly or
indirectly subset \texttt{Objects::Object::ownedPerformances} (\ref{ownedPerformances}).

\ksec{Functions Semantics (KerML 8.4.4.8)}\label{sec:functions}

\pn Functions\index{function} are not behaviors.  They merely determine values.
This is a significant change for SFS requiring a change in the specialization of metaclasses (\ref{fig:classdiagram}).

\pn The \texttt{Function} metaclass specializes \texttt{Classifier} instead of \texttt{Behavior}.
\ecaption{metaclass Function}
\begin{lstlisting}
  metaclass Function specializes Classifier {
  	derived var feature expression : Expression[0..*] subsets 'feature';
  	derived var feature result : Feature[1..1] subsets output;
  }		
\end{lstlisting}

\pn A \texttt{Function} directly, indirectly, or implicitly specializes \texttt{Performances::Evaluation} which specializes \texttt{Anything}
instead of \texttt{Performance}.

\ecaption{Performances::Evaluation}
\begin{lstlisting}
	abstract function Evaluation specializes Anything {
		doc
		/*
		 * Evaluation is the most general class of functions that may be evaluated to determine
		 * a result.
		 */
		return result: Anything[0..*] nonunique;
	}
\end{lstlisting}

\pn To model computation by the system use  \texttt{Performances::Computation}. 
\ecaption{Performances::Computation}
\begin{lstlisting}
	abstract behavior Computation specializes Performance, Evaluation {
		doc
		/*
		 * Computation is a specialization of Performance and Evaluation for the case of
		 * computing a result.
		 */
			
	}
\end{lstlisting}

\pn The feature \texttt{Performances::evaluations} is changed to subset  \texttt{things} rather than  \texttt{performances}. 
\ecaption{Performances::evaluations}
\begin{lstlisting}
	abstract expr evaluations: Evaluation[0..*] nonunique subsets things {
		doc
		/*
		 * evaluations is a specialization of things for evaluations of Functions.
		 */
	}
\end{lstlisting}

\pn Because behavior \texttt{Computation} was added, it needs a corresponding feature, \texttt{computations}.
\begin{lstlisting}
	abstract step computations: Computation[0..*] nonunique subsets performances {
		doc
		/*
		 * computations is a specialization of performances for the Computation.
		 */
	}
\end{lstlisting}


\pn The value of the result can be specified as a naked expression at the end:

\ecaption{function}
\begin{lstlisting}
function Average {
  in scores[1..*] : Rational; 
  return : Rational;
  sum(scores) / size(scores)
}
\end{lstlisting}

or bound to the return parameter:
\ecaption{bound return parameter}
\begin{lstlisting}
function Average {
  in scores[1..*] : Rational; 
  return : Rational = sum(scores) / size(scores);
}
\end{lstlisting}

%No special timing semantics for \texttt{Function} other than being a \texttt{Behavior}.
%
%However, in SFS, \texttt{Expression} is a \texttt{Value} not a \texttt{Step}.  If an expression is to be evaluated by the system being modeled, wrap it in a function to be a \texttt{Performances::Evaluation} so it is an \texttt{Occurrence}.

